\section{Methods}
\label{sec:methods}

\subsection{Splits}

For each pool we draw a single random partition into training, validation and
test sets (\SI{90}{\percent}/\SI{5}{\percent}/\SI{5}{\percent}) after removing
failed docking runs. The same partition is reused by every backbone, so a
comparison between representations is never confounded by a difference in the
underlying rows. Training indices are stored in random order, which allows a
training budget of $N$ to be realised as the first $N$ indices; learning-curve
points are therefore nested subsets rather than independent samples, and the
corresponding embedding cache for budget $N$ is a prefix of the full cache.

Model selection uses the validation set only. The test set is read once per
run, after the best epoch has been selected, and the resulting predictions are
written to disk. This last detail is what makes the fusion and size-control
analyses in Sections~\ref{sec:results-fusion} and~\ref{sec:results-size}
possible without retraining.

\subsection{Prediction head}

To isolate the effect of the representation, the prediction head is held fixed
across all backbones: a multilayer perceptron with hidden widths
$(1024, 512)$, batch normalisation, ReLU activations and dropout $0.1$,
regressing a standardised score under mean squared error. Inputs are
standardised per feature using statistics estimated on a training subsample.
We optimise with AdamW (learning rate \num{e-3}, weight decay \num{e-5}) under
a cosine schedule, for at most \num{40} epochs with early stopping after
\num{6} epochs without validation improvement.

Holding the head fixed trades per-backbone optimality for comparability. We
address the resulting objection --- that a single architecture may suit some
representations better than others --- in Section~\ref{sec:results-robust}.

\subsection{Descriptor baseline}
\label{sec:methods-baseline}

The baseline replaces the learned embedding with \num{12} physicochemical
descriptors computed with RDKit~\citep{rdkit}: heavy atom count, molecular
weight, \mbox{cLogP}, topological polar surface area, ring count, rotatable
bond count, hydrogen bond donor and acceptor counts, fraction of sp\textsuperscript{3}
carbons, halogen count, aromatic ring count, and heteroatom count. These pass
through the identical split, prediction head, optimiser and metrics as every
learned representation, so the comparison is of representation alone.

\subsection{Metrics}

Lower scores indicate better predicted binders, so ``top-$k$'' refers to the
$k$ most negative values. We report regression metrics (RMSE, MAE, $R^2$,
Pearson $r$, Spearman $\rho$) and screening metrics: recall of the true top
fraction $f$ within the predicted top fraction $f$, and the corresponding
enrichment factor $\mathrm{EF}@f = \mathrm{recall}@f / f$, at
$f \in \{0.001, 0.01, 0.05\}$.

The two families answer different questions and need not agree. Pearson $r$
and RMSE are computed over the whole pool and are therefore dominated by the
large majority of non-binders; enrichment at \SI{0.1}{\percent} describes the
only region of the ranking a screening campaign would act upon.

\subsection{Size-controlled evaluation}
\label{sec:methods-size}

Because docking scores scale with molecular size, we compute three additional
quantities for every set of predictions, all from the stored test predictions
and therefore applicable retrospectively.

\paragraph{Partial correlation.}
The correlation between true and predicted score with heavy-atom count $s$
partialled out of both,
\begin{equation}
r_{yp \cdot s} = \frac{r_{yp} - r_{ys} r_{ps}}
                      {\sqrt{(1 - r_{ys}^2)(1 - r_{ps}^2)}},
\end{equation}
which removes a linear dependence on size.

\paragraph{Residual correlation.}
The correlation between the residuals of $y$ and of $\hat{y}$ after each is
regressed on a degree-3 polynomial in $s$, which permits size to act
non-linearly.

\paragraph{Within-stratum metrics.}
The test set is partitioned into ten equal-count bins of heavy-atom count, and
Pearson $r$, Spearman $\rho$ and $\mathrm{EF}@0.01$ are computed inside each
bin and averaged with weights proportional to bin size. This assumes no
functional form for the size dependence, and corresponds to the practical
question of whether a model can rank molecules of comparable size.

As a reference point we also evaluate the degenerate predictor
$\hat{y} = -s$, which by construction attains zero partial and residual
correlation and near-chance within-stratum performance. Its raw correlation
quantifies how much of a pool's score variance size alone explains.

\subsection{Fusion}

We combine representations by late fusion: each backbone is trained
independently and their test predictions are averaged. This permits any subset
to be evaluated without retraining, which is the property that made late fusion
attractive in prior multimodal
benchmarks~\citep{martinezgalindo2026multimodality}. We additionally consider
rank averaging and standardised averaging, which discard calibration in
exchange for robustness to differences in scale between members.
\todo{Add intermediate fusion (concatenated embeddings, jointly trained) if the
multimodality framing is retained.}
